\section{Institutional Ownership and Crowded Trades}
\label{sec:crowded_trades}

Institutional investors, unlike individual investors, typically manage large portfolios and may therefore generate economically meaningful demand for individual securities. Their investment decisions can affect prices not only through the information they incorporate into their trades, but also through the size and similarity of their portfolio allocations. Early evidence on institutional demand shows that institutions display systematic preferences for particular stock characteristics and that the growth and composition of institutional ownership can affect relative stock prices and returns \parencite{gompers2001institutional}. This makes the composition of the institutional investor base potentially relevant for understanding both price formation and the distribution of risk across securities.

A related strand of the literature examines whether institutional investors trade in a correlated manner. \textcite{wermers1999mutual} studies mutual fund trading and finds evidence of herding, particularly among smaller stocks and growth-oriented funds, although average levels of herding are relatively modest. \textcite{sias2004institutional} further documents persistence in institutional demand and provides evidence consistent with institutions following one another into and out of the same securities. These findings suggest that institutional portfolio decisions cannot necessarily be viewed as independent: common information, persistence in institutions' own trading, and inference from other investors' trades can generate correlated trading behavior across institutions.

It is useful, however, to distinguish \emph{herding} from \emph{crowding}. Herding generally refers to the trading process through which investors buy or sell the same securities over a given period. Crowding instead describes the resulting state in which many investors are exposed to the same security, strategy, or set of positions at a particular point in time. A position may therefore be crowded even if the investors currently holding it are not actively trading in the same direction \parencite{wermers1999mutual,brown2022crowded}. This distinction is important for the present thesis, which focuses primarily on the structure of observed institutional holdings rather than on directly identifying simultaneous trades in real time. Figure~\ref{fig:herding-vs-crowding} illustrates this conceptual distinction.

\begin{figure}[H]
    \centering
    \includegraphics[width=0.9\textwidth]{08_04_herding_vs_crowding_networkx.pdf}
    \caption{Conceptual distinction between herding and crowding in institutional ownership networks.}
    \label{fig:herding-vs-crowding}
\end{figure}

\noindent
From a risk perspective, crowding becomes particularly relevant when the investors sharing a position are themselves similar. Aggregate institutional ownership alone does not reveal whether a stock is held by a diverse group of investors or by institutions with strongly overlapping portfolios. In the latter case, common shocks to investor capital, investment strategies, or risk constraints may generate correlated demand for liquidity. The ownership structure may consequently create an additional channel through which stress affecting investors can be transmitted to the securities they jointly hold. This intuition connects the literature on institutional ownership to the broader literature on flow-driven price pressure and fire sales, discussed in Section~\ref{sec:liquidity_fire_sales}.

More recent work has examined crowdedness directly as a security-level characteristic. \textcite{brown2022crowded} use hedge-fund position data to construct measures of security-level crowdedness and show that greater exposure to crowdedness is associated with downside tail risk and larger drawdowns during periods of market distress. Their results are important for the present analysis because they suggest that information contained in institutional portfolios may be relevant not only for expected returns but also for asymmetric downside outcomes. Related evidence from \textcite{lou2022comomentum} shows that excessive commonality in arbitrage activity can be associated with subsequent reversals, supporting the broader idea that crowded investment strategies may become destabilizing when many investors pursue similar positions.
